\documentclass[10pt]{article}
\usepackage[margin=.85in]{geometry}
\usepackage{amsmath,amssymb}
\usepackage[hidelinks]{hyperref}

\title{Conjecture 00000000582: Bergman \texorpdfstring{$h$}{h}-vectors}
\author{}
\date{}

\begin{document}
\maketitle

Take the rational vector matroids $M,N$ whose columns are the lifts
$(x_i,y_i,1)$ of the following six labeled points in each row.
\[
\begin{array}{c|rrrrrr}
 &1&2&3&4&5&6\\ \hline
 M&(0,0)&(1,0)&(2,0)&(3,0)&(0,1)&(5,2)\\
 N&(0,0)&(2,0)&(0,2)&(1,0)&(1,1)&(0,1)
\end{array}
\]
These are simple rank-three matroids. Their points are distinct,
and the triples $125$ for $M$ and $123$ for $N$ are noncollinear.
Every rank-two flat is the closure of a pair $i,j$. Its elements $k$
are exactly those satisfying
$(x_j-x_i)(y_k-y_i)-(y_j-y_i)(x_k-x_i)=0$.
Testing the fifteen pairs and removing repeated closures gives
\[
\begin{array}{c|l}
M &1234,\ 15,\ 16,\ 25,\ 26,\ 35,\ 36,\ 45,\ 46,\ 56\\
N &124,\ 136,\ 235,\ 15,\ 26,\ 34,\ 45,\ 46,\ 56
\end{array}
\]
where a string denotes the set of its labels. For either matroid,
the M\"obius value from the empty flat to a point is $-1$, and the
value for a line $L$ is $|L|-1$ by the M\"obius recurrence.
The sum over lines is $3+9=12$ in $M$ and $3\cdot2+6=12$ in $N$.
The top flat therefore has value $-(1-6+12)=-7$ in both, giving
$\chi(t)=t^3-6t^2+12t-7$.

By Ardila and Klivans \cite[Theorem~1]{ak}, the fine Bergman complex
is the order complex of the nonempty proper flats. Its vertices are
the points and lines, and its edges are the incidences $i\in L$.
Thus $(f_0,f_1)=(16,22)$ for $M$ and $(15,21)$ for $N$. The
one-dimensional face-to-$h$ transform is
\[
(h_0,h_1,h_2)=(1,f_0-2,f_1-f_0+1).
\]
This yields $(1,14,7)$ for $M$ and $(1,13,7)$ for $N$.

The coarse complexes are also one-dimensional, and their $h$-vectors
differ as well.
Both matroids are connected, since a disconnected simple rank-three
matroid has a rank-one component consisting of a coloop, and neither
list has a line through five points. By Feichtner and Sturmfels
\cite[Proposition~2.6 and Theorem~2.7]{fs}, coarse vertices correspond
to proper flats whose restrictions and contractions are connected.
Every point is such a flat. Its contraction is loopless of rank two,
with parallel classes corresponding to the lines through the point.
There are at least three such lines at every point in the lists, while
a disconnected loopless rank-two matroid has exactly two parallel
classes. For a line $L$, the contraction is a loopless rank-one
matroid, and the restriction is the uniform matroid of rank two on
$|L|$ points. This restriction is connected exactly when $|L|\ge3$.
The coarse vertex counts are consequently $6+1=7$ and $6+3=9$,
so the coarse values of $h_1=f_0-2$ differ as well.

The shared characteristic polynomial therefore cannot determine the
Bergman $h$-vector, including through a fixed Stirling-type matrix.

The Lean certificate enumerates all pair closures by the determinant,
checks the rank and incidence conditions, and derives the equal
characteristic coefficients and unequal $h$-data. The identifications
with the fine and coarse Bergman complexes use the cited theorems.

\begin{thebibliography}{9}
\bibitem{ak} F. Ardila and C. J. Klivans,
\emph{The Bergman complex of a matroid and phylogenetic trees},
J. Combin. Theory Ser. B 96 (2006), 38--49.
\href{https://arxiv.org/abs/math/0311370}{arXiv:math/0311370}.
\bibitem{fs} E. M. Feichtner and B. Sturmfels,
\emph{Matroid polytopes, nested sets and Bergman fans},
Portugaliae Math. 62 (2005), 437--468.
\href{https://arxiv.org/abs/math/0411260}{arXiv:math/0411260}.
\end{thebibliography}

\end{document}
